\fancyhead[L]{\small BESSEL ÜBER DEN KOMETEN VON 1807.}
\hypertarget{native-535lower}{}\label{native-535lower}
\begin{center}\small Göttingische gelehrte Anzeigen. Stück 26. S. 250..256. 16. Februar 1811.\\[3pt]\rule{.72\linewidth}{.3pt}\end{center}

\emph{Untersuchungen über die scheinbare und wahre Bahn des im Jahr \textup{1807} erschienenen großen Kometen, von} F. W. \textsc{Bessel}, Professor der Astronomie in Königsberg. Bei \textsc{Friedr. Nicolovius}. Königsberg. 82 Seiten im Quartformat.

Der Komet von 1807, heller als jeder seit vielen Jahren sichtbare Komet, galt bekanntlich als eine der bemerkenswertesten Himmelserscheinungen. Für die Astronomen war er wegen seiner langen Sichtbarkeit noch interessanter. Die Beobachter nutzten diese Zeit gleichsam im Wetteifer, um seine scheinbare Bewegung möglichst vollständig und genau zu bestimmen. Auch die Berechnung seiner parabolischen Bahn beschäftigte mehrere Astronomen. Niemand hat jedoch so viel Sorgfalt auf die Untersuchung der Theorie dieses Kometen verwendet wie der Verfasser der vorliegenden Schrift, der sich bekanntlich schon um mehrere ältere Kometen, besonders um den von 1769, sehr verdient gemacht hatte. Ein großer Teil der Untersuchungen des Verfassers über den Kometen von 1807 war bereits durch die Monatliche Crrespondenz und das Astronomische Jahrbuch bekannt geworden. Die vorliegende Schrift soll all dies im Zusammenhang darstellen und zugleich die späteren Arbeiten vorlegen, mit denen er die Theorie dieses Kometen zum Abschluss brachte. Die Schrift liefert einen neuen Beleg sowohl für den heutigen Stand der Beobachtungskunst als auch für die Leistungsfähigkeit der Rechnung, wenn eine so geschickte Hand sie leitet.

Der \emph{erste Abschnitt} behandelt die scheinbare Bahn des Kometen am Himmel. Herr \textsc{Bessel} wollte nicht alle Beobachtungen des Kometen sammeln, sondern nur so viele sorgfältig angestellte Beobachtungen zusammenführen, wie für eine hinreichende Darstellung der scheinbaren Bewegung während seiner gesamten Sichtbarkeit vom 22. September 1807 bis zum 27. März 1808 nötig waren. Er teilt daher zunächst seine eigenen, in Lilienthal mit dem Ringmikrometer angestellten Beobachtungen mit. Angesichts der gewissenhaften Sorgfalt und der darin hervorragenden Geschicklichkeit des Verfassers verdienen sie besonderes Vertrauen. Die verschiedenen Vorsichtsregeln und Bemerkungen zu dieser Art von Beobachtungen, die er bei dieser Gelegenheit mitteilt, werden vielen Lesern sehr willkommen sein. Den Durchmesser des Gesichtsfelds bestimmte \textsc{Bessel} anhand von Sternpaaren, deren Deklinationen genau bekannt sind und sich fast um diesen Durchmesser unterscheiden. Dieses Verfahren ist sehr zu empfehlen, und man kann dann, wie Bessel es tut, den Durchmesser des Felds durch eine einfache Reihe darstellen. Sobald der Durchmesser den Deklinationsunterschied der Sterne jedoch um mehr als eine Bogenminute übertrifft, verwendet der Rezensent lieber die Formeln \[
\frac{a'-a}{2d}=\mathrm{tang}\,\varphi,\qquad \frac{a'+a}{2d}=\mathrm{tang}\,\psi,\qquad \frac{d}{\cos\varphi\cos\psi}=D
\] wobei $a$ und $a'$ die beiden Sehnen, $d$ den Deklinationsunterschied der Sterne und $D$ den Durchmesser des Gesichtsfelds bezeichnen. Die eigenen Beobachtungen von \textsc{Bessel} reichen vom 4. Oktober 1807 bis zum 24. Februar \SourceLeaf{536}1808. Bei ihrer Mitteilung hat er in sehr lobenswerter Weise so viele Einzelheiten beigefügt, dass man die Beobachtungen künftig erneut reduzieren kann, wenn die Vergleichssterne, die größtenteils der Histoire céleste entnommen sind, genauer bekannt sind. Darauf folgen die Beobachtungen von Dr. \textsc{Olbers} in Bremen vom 8. Oktober 1807 bis zum 14. Februar 1808, die Beobachtungen von \textsc{Thulis} in Marseille vom 22. September bis zum 2. Oktober und die von \textsc{Wisniewsky} in Petersburg vom 18. bis zum 27. März. Die beiden letztgenannten Reihen hat der Verfasser neu reduziert. Schließlich folgen, um die im Februar 1808 noch verbliebene Lücke zu füllen, die Beobachtungen von \textsc{Oriani} in Mailand vom 13. bis zum 28. Februar.

Der \emph{zweite Abschnitt} ist der Bestimmung der wahren Bahn des Kometen gewidmet und gliedert sich seinerseits in drei Abteilungen. Die \emph{erste Abteilung} gibt parabolische und rein elliptische Elemente des Kometen an. Die ersten parabolischen Elemente berechnete der Verfasser schon am 23. Oktober. Die weiteren Beobachtungen erlaubten bereits am 6. November eine Verbesserung. Im März 1808 berechnete er ein drittes System parabolischer Elemente, das sich den ersten und letzten ihm damals bekannten Beobachtungen und den Rektaszensionen der dazwischenliegenden Beobachtungen möglichst genau anschloss. Bei den dazwischenliegenden Deklinationen blieb jedoch noch ein Unterschied von einer Bogenminute. Er hätte sich zwar noch zu einem erheblichen Teil auf die fünf übrigen Größen verteilen lassen, wäre aber dennoch zu groß geblieben, als dass man daraus nicht auf eine merkliche Abweichung von der Parabel hätte schließen müssen. \textsc{Bessel} gab daher die parabolische Hypothese auf und bestimmte unabhängig von ihr ein viertes System von Elementen. Diese ergeben eine Ellipse mit einer Umlaufzeit von 1953 Jahren. Im Vergleich mit allen Beobachtungen aus Lilienthal und Bremen zeigten sich überall nur so unregelmäßig verlaufende Unterschiede, dass sie allein den Fehlern der einzelnen Beobachtungen zugeschrieben werden konnten. Bald darauf zeigten jedoch die ersten Beobachtungen aus Marseille, dass jene Elemente die Rektaszensionen etwas zu groß ergaben. Das veranlasste ihn zur Berechnung eines fünften Systems von Elementen, wodurch sich die Umlaufzeit auf 1483 Jahre verringerte. Diese Elemente stellten alle 91 Beobachtungen so gut dar, dass sich nirgends mehr mit Sicherheit ein eindeutiger Unterschied angeben ließ. Damit war die Bahn so genau bestimmt, wie es mit rein elliptischen Elementen möglich war.

Es war jedoch zu erwarten, dass die Berücksichtigung der Störungen, die der Komet während seiner Sichtbarkeit durch die Planeten unseres Sonnensystems erfahren musste, diese Ergebnisse umso stärker verändern würde, weil die Umlaufzeit so stark von der kleinen Differenz zwischen der Exzentrizität und Eins abhängt. Selbst eine an sich kleine Änderung dieser Differenz kann die Umlaufzeit gewaltig verändern. Es war daher eine interessante Aufgabe, diese Störungen angemessen in die Rechnung einzubeziehen, da sich sonst überhaupt nicht beurteilen ließ, innerhalb welcher Grenzen man die Umlaufzeit ansetzen durfte. Diesem Zweck dienen die beiden folgenden Abteilungen. In der \emph{zweiten Abteilung} entwickelt der Verfasser nämlich die allgemeine Methode zur Berechnung der Störungen des Kometen. Bekanntlich behandelt man dabei am zweckmäßigsten die Bahn des Kometen als veränderliche Ellipse, berechnet die täglichen Änderungen ihrer einzelnen Bestimmungsgrößen durch die störenden Kräfte numerisch und bestimmt deren Gesamtwirkung durch mechanische Quadraturen. Die Formeln für die augenblickliche Änderung der Elemente einer elliptischen Bahn durch störende Kräfte sind zwar schon von mehreren Geometern hergeleitet worden. Dennoch wird man die neue Herleitung, die \textsc{Bessel} in dieser Abteilung mit großer analytischer Eleganz gibt, umso lieber lesen, als bei einer Ellipse von so großer Exzentrizität zum Teil eine Abänderung nötig war. Statt der Epoche der mittleren Länge oder Anomalie wurde nämlich die Zeit des Durchgangs durch das Perihel gewählt. Der sich leicht ergebende \SourceLeaf{537}Ausdruck für die augenblickliche Störung dieser Durchgangszeit ließ sich wegen des genannten Umstands nicht gut zur Rechnung verwenden und musste erst eigens umgeformt werden.

In der \emph{dritten Abteilung} wendet der Verfasser diese Methode schließlich auf den Kometen von 1807 an, indem er zunächst die Summe der störenden Kräfte für sieben gleich weit auseinanderliegende Zeitpunkte vom 22. September 1807 bis zum 20. März 1808 berechnet. Er sagt nicht ausdrücklich, ob er alle Planeten oder wie viele von ihnen er dabei berücksichtigt hat. Der Rezensent hätte den Beitrag jedes einzelnen Planeten gern gesondert angegeben gefunden, da ein Überblick über dieses Verhältnis nicht uninteressant gewesen wäre. Darauf berechnet er für jene sieben Zeitpunkte den Betrag der täglichen Änderung jedes der sechs Elemente. Er betrachtete nun das zuletzt gefundene System von Elementen als für den 22. September 1807 gültig und berechnete, nachdem er ihre augenblicklichen Änderungen mechanisch integriert hatte, ihre Werte für sechs verschiedene Zeitpunkte vom 28. September 1807 bis zum 23. März 1808. Für diese Zeitpunkte hatte er zugleich die geozentrischen Orte jeweils aus allen benachbarten beobachteten Orten mit größter Sorgfalt bestimmt. Der Vergleich dieser Orte mit der Rechnung nach jenen veränderlichen Elementen ergab zusammen mit den numerisch hergeleiteten Differentialbeziehungen zwischen den einzelnen Elementen und den berechneten geozentrischen Längen und Breiten 12 lineare Gleichungen für die noch zu bestimmenden Korrekturen der Elemente. Daraus wurden die Werte dieser Korrekturen unter gebührender Berücksichtigung des relativen Gewichts der einzelnen Positionen nach der Methode der kleinsten Quadrate abgeleitet. So ergaben sich schließlich die endgültigen Elemente für den 22. September, von denen wir hier nur die große Halbachse $= 143{,}195$ und die Umlaufzeit $= 1713\tfrac{1}{2}$ Jahre anführen. Diese Umlaufzeit wird jedoch durch die fortwirkenden Störungen ständig verändert. Selbst wenn dieses Ergebnis im Übrigen vollkommen zuverlässig wäre, ließe sich der Zeitpunkt der tatsächlichen Wiederkehr daher erst nach einer endlos umfangreichen Rechnung angeben. Aus einer überschlägigen Berechnung der von Jupiter verursachten Störungen in den Jahren 1808..1815, gegenüber denen die Störungen durch die anderen Planeten unmerklich werden, findet \textsc{Bessel} bereits eine Verkürzung der Umlaufzeit um 170 Jahre.

Zum Schluss untersucht der Verfasser noch, welchen Grad an Genauigkeit man diesen Ergebnissen zuschreiben darf. Nimmt man für Länge und Breite bei den ersten drei Normalorten am 28. September, 22. Oktober und 11. November einen Fehler von $5''$ an, bei den Orten vom 8. Dezember und 21. Februar einen doppelt so großen und beim Ort vom 23. März einen viermal so großen, so könnte sich die Umlaufzeit bei der ungünstigsten Kombination dieser Fehler auf 2157 Jahre erhöhen oder auf 1404 Jahre verringern. Eine solche Kombination ist zwar möglich, aber höchst unwahrscheinlich. Da \textsc{Bessel} so zahlreiche Beobachtungen mit solcher Sorgfalt zur Bestimmung der Normalorte verwendet hat, ist auch kaum zu befürchten, dass darin noch größere Fehler als die eben angenommenen zurückgeblieben sind. Man darf also wohl hoffen, dass die vorhin bestimmte Umlaufzeit auf hundert Jahre genau ist und der Komet ungefähr im vierunddreißigsten Jahrhundert unserer Zeitrechnung wiederkehrt. Unsere späten Nachkommen werden sich dann dankbar an die Arbeit von \textsc{Bessel} erinnern.

Diese letzten Überlegungen führen zugleich auf einfache Weise zu der vollkommenen Überzeugung, dass der Komet \emph{gewiss} eine Ellipse und weder eine Parabel noch eine Hyperbel beschreibt. Um eine nichtelliptische Bahn für möglich zu halten, müsste man bei den Normalorten \emph{mindestens} siebenmal so große Fehler zulassen. Das kann man ohne Bedenken für schlechthin unmöglich erklären. So gehört also (bei der heutigen Vollkommenheit der \emph{praktischen Astronomie}) der Komet von 1807 zu den wenigen, von denen wir mit Gewissheit behaupten können, dass sie dereinst gewiss wiederkehren werden, wenn keine anderen als die uns bekannten Kräfte auf sie einwirken.
